% Tugas 4 STSI4310 Metodologi Penelitian
% Kompilasi: pdflatex Tugas4.tex (jalankan dua kali agar nomor halaman daftar isi terisi)
\documentclass[12pt,a4paper]{article}
\usepackage[a4paper,top=3cm,bottom=3cm,left=4cm,right=3cm]{geometry}
\usepackage[T1]{fontenc}
\usepackage[utf8]{inputenc}
\usepackage{mathptmx} % Times
\usepackage{courier} % font monospace Type 1, hindari bitmap EC
\usepackage[protrusion=true,expansion=true]{microtype} % kurangi jarak antarkata yang renggang
\usepackage{setspace}
\usepackage{graphicx}
\usepackage{enumitem}
\usepackage{caption}
\usepackage{array,tabularx,booktabs}
\usepackage{url}
\usepackage[hidelinks,unicode]{hyperref}
\urlstyle{same}

% Spasi 1,5, paragraf menjorok, tanpa jarak antarparagraf
\onehalfspacing
\setlength{\parindent}{1.25cm}
\setlength{\parskip}{0pt}

% Tanpa pemenggalan kata: pola pemenggalan bawaan berbahasa Inggris
% menghasilkan pemenggalan yang salah untuk bahasa Indonesia.
\hyphenpenalty=10000
\exhyphenpenalty=10000
\setlength{\emergencystretch}{3em}

% Penanda istilah asing (dicetak miring secara konsisten)
\newcommand{\asing}[1]{\textit{#1}}
\newcommand{\pim}{\textit{pi-memory}}
\newcommand{\phm}{\textit{pi-hermes-memory}}

% Judul bertingkat: BAB (tengah), subbab (x.y), sub-subbab (x.y.z)
\newcommand{\bab}[3]{%
  \clearpage\phantomsection\label{#3}%
  \begin{center}\large\bfseries BAB #1\\[0.2em] #2\end{center}\vspace{0.8em}}
\newcommand{\subbab}[2]{%
  \par\vspace{0.8em}\phantomsection\label{#2}%
  \noindent\textbf{#1}\par\nopagebreak\vspace{0.3em}}
\newcommand{\subsubbab}[2]{%
  \par\vspace{0.5em}\phantomsection\label{#2}%
  \noindent\textbf{#1}\par\nopagebreak\vspace{0.2em}}

% Entri daftar isi: tingkat 0 (bab), 1 (subbab), 2 (sub-subbab)
\newcommand{\tocentry}[3]{%
  \ifcase#1\relax
    \vspace{0.5em}\noindent\textbf{#2}\dotfill\textbf{\pageref{#3}}\par
  \or
    \noindent\hspace*{1.25cm}#2\dotfill\pageref{#3}\par
  \or
    \noindent\hspace*{2.25cm}#2\dotfill\pageref{#3}\par
  \fi}

\begin{document}

% ------------------------------------------------------------------
% HALAMAN COVER (sesuai Lampiran Panduan Penyusunan Proposal STSI4310)
% ------------------------------------------------------------------
\begin{titlepage}
\thispagestyle{empty}
\centering
\begin{singlespace}
{\large\bfseries Tugas 4}\par
\vspace{0.4em}
{\large\bfseries STSI4310 -- Metodologi Penelitian}\par

\vfill

{\large\bfseries Akurasi dan Efisiensi Penyimpanan Memori Agen Pi\\
melalui Evaluasi Komparatif\\
Gate Klasifikasi Terkalibrasi}\par

\vfill

\includegraphics[width=5.2cm]{media/logo-universitas-terbuka.png}\par

\vfill

{\large\bfseries Fadhil Andriawan - 053497355}\par

\vfill

{\large\bfseries Program Studi Sistem Informasi}\par
\vspace{0.3em}
{\large\bfseries Fakultas Sains dan Teknologi}\par
\vspace{0.3em}
{\large\bfseries Universitas Terbuka}\par
\vspace{0.3em}
{\large\bfseries 2026}\par
\end{singlespace}
\end{titlepage}

% ------------------------------------------------------------------
% DAFTAR ISI
% ------------------------------------------------------------------
\pagenumbering{roman}
\phantomsection\label{sec:daftar-isi}
\begin{center}\large\bfseries DAFTAR ISI\end{center}
\vspace{0.8em}
\begin{singlespace}
\tocentry{0}{DAFTAR ISI}{sec:daftar-isi}
\tocentry{0}{BAB I PENDAHULUAN}{sec:bab1}
\tocentry{1}{1.1 Latar Belakang Masalah}{sec:1.1}
\tocentry{1}{1.2 Rumusan Masalah}{sec:1.2}
\tocentry{1}{1.3 Tujuan dan Manfaat Penelitian}{sec:1.3}
\tocentry{2}{1.3.1 Tujuan Penelitian}{sec:1.3.1}
\tocentry{2}{1.3.2 Manfaat Penelitian}{sec:1.3.2}
\tocentry{0}{BAB II KAJIAN PUSTAKA DAN LANDASAN TEORI}{sec:bab2}
\tocentry{1}{2.1 Kajian Pustaka dan Landasan Teori}{sec:2.1}
\tocentry{2}{2.1.1 Memori Lintas Sesi pada \asing{Coding Agent}}{sec:2.1.1}
\tocentry{2}{2.1.2 Penilaian Kelayakan dan Kendali Jalur Tulis Memori}{sec:2.1.2}
\tocentry{2}{2.1.3 Klasifikasi Kelayakan Penyimpanan}{sec:2.1.3}
\tocentry{2}{2.1.4 Kalibrasi Probabilitas \asing{Gate}}{sec:2.1.4}
\tocentry{2}{2.1.5 \asing{Band} Ketidakpastian dan Pemeriksaan LLM}{sec:2.1.5}
\tocentry{2}{2.1.6 Evaluasi Efisiensi dan Mutu Memori}{sec:2.1.6}
\tocentry{2}{2.1.7 Posisi Penelitian}{sec:2.1.7}
\tocentry{1}{2.2 Pertanyaan Penelitian dan Hipotesis}{sec:2.2}
\tocentry{0}{BAB III METODE PENELITIAN}{sec:bab3}
\tocentry{1}{3.1 Metode Penelitian}{sec:3.1}
\tocentry{2}{3.1.1 Pendekatan dan Jenis Penelitian}{sec:3.1.1}
\tocentry{2}{3.1.2 Objek, Populasi, dan Sampel Penelitian}{sec:3.1.2}
\tocentry{2}{3.1.3 Variabel Penelitian}{sec:3.1.3}
\tocentry{2}{3.1.4 Desain dan Prosedur Eksperimen}{sec:3.1.4}
\tocentry{1}{3.2 Teknik Pengumpulan Data dan Instrumen Penelitian}{sec:3.2}
\tocentry{2}{3.2.1 Teknik Pengumpulan Data}{sec:3.2.1}
\tocentry{2}{3.2.2 Instrumen Penelitian}{sec:3.2.2}
\tocentry{2}{3.2.3 Validitas dan Reliabilitas Instrumen}{sec:3.2.3}
\tocentry{0}{BAB IV ANALISIS DATA DAN KESIMPULAN SEMENTARA}{sec:bab4}
\tocentry{1}{4.1 Teknik Analisis Data Penelitian}{sec:4.1}
\tocentry{2}{4.1.1 Statistika Deskriptif}{sec:4.1.1}
\tocentry{2}{4.1.2 Statistika Inferensial dan Uji Hipotesis}{sec:4.1.2}
\tocentry{1}{4.2 Kesimpulan Sementara dan Batasan Penelitian}{sec:4.2}
\tocentry{2}{4.2.1 Kesimpulan Sementara}{sec:4.2.1}
\tocentry{2}{4.2.2 Batasan Penelitian}{sec:4.2.2}
\tocentry{0}{DAFTAR PUSTAKA}{sec:daftar-pustaka}
\end{singlespace}

\clearpage
\pagenumbering{arabic}

% ------------------------------------------------------------------
% BAB I
% ------------------------------------------------------------------
\bab{I}{PENDAHULUAN}{sec:bab1}

\subbab{1.1 Latar Belakang Masalah}{sec:1.1}

\asing{Coding agent} seperti Pi bekerja dalam sesi. Ketika sesi berakhir, preferensi pengguna, konvensi repositori, keputusan arsitektur, dan pekerjaan yang belum selesai tidak selalu tersedia pada sesi berikutnya. Akibatnya, pengguna perlu menjelaskan kembali informasi yang sama dan agen berpotensi mengulang kesalahan. \asing{Extension} \pim{} dan \phm{} menyediakan memori lintas sesi untuk mengatasi masalah ini (chandra447, 2026; jayzeng, 2026). Pada \phm{}, \asing{review} latar belakang dilakukan secara bawaan setiap 10 \asing{turn} atau 15 \asing{tool call} menggunakan panggilan model generatif (chandra447, 2026), sehingga proses penilaian dapat berulang selama sesi berlangsung.

Penyatuan penilaian kelayakan dan penulisan entri menimbulkan persoalan efisiensi dan keandalan keputusan. Penilaian melalui model generatif tetap menggunakan token konteks meskipun kandidat akhirnya ditolak. Jika model menggunakan \asing{sampling}, keputusan untuk masukan yang sama juga dapat berbeda. Selain itu, nilai keyakinan model belum tentu sesuai dengan proporsi keputusan yang benar (Guo et al., 2017).

Beberapa pendekatan telah memisahkan penilaian dari penulisan. Command Code menggunakan model neuro-simbolik untuk menilai sinyal preferensi dan model generatif untuk menulis hasilnya (Command Code, 2026). SAGE memisahkan kandidat yang jelas baru, berulang, dan ambigu, lalu hanya meneruskan kandidat ambigu ke LLM (Wang et al., 2026). Penghematan biaya dan mutu keputusan dari pemisahan ini masih perlu diuji pada memori sesi \asing{coding agent}, yang memuat keputusan teknis, konvensi repositori, dan koreksi pengguna. Data percakapan jangka panjang yang digunakan dalam penelitian sebelumnya juga memiliki konteks yang berbeda dengan sesi pemrograman (Maharana et al., 2024).

Berdasarkan permasalahan tersebut, penelitian ini membandingkan \asing{gate} klasifikasi terkalibrasi dengan mekanisme penyimpanan pada \pim{} dan \phm{}. Kandidat yang probabilitasnya berada dalam \asing{band} ketidakpastian diteruskan ke LLM, sedangkan kandidat lainnya diputuskan secara otomatis oleh \asing{gate}. Model generatif tetap digunakan untuk menulis entri yang diterima. Evaluasi mencakup akurasi keputusan, kalibrasi, penggunaan token, biaya, latensi, serta mutu memori pada sesi lanjutan.

Penelitian ini berada dalam ruang lingkup Sistem Informasi karena membahas perancangan dan evaluasi komponen sistem berbasis kecerdasan artifisial serta pengaruhnya terhadap biaya operasional dan mutu informasi yang disimpan.

\subbab{1.2 Rumusan Masalah}{sec:1.2}

Rumusan masalah penelitian ini disusun dalam bentuk pertanyaan penelitian sebagai berikut.
\begin{enumerate}[leftmargin=1.25cm]
\item Bagaimana akurasi keputusan \asing{gate} klasifikasi terkalibrasi dibandingkan dengan mekanisme penyimpanan pada \pim{} dan \phm{}?
\item Apakah probabilitas yang dihasilkan \asing{gate} terkalibrasi dan dapat digunakan untuk menentukan keputusan?
\item Bagaimana perbandingan penggunaan token, biaya, dan latensi per operasi tulis antara mekanisme yang diusulkan, \pim{}, dan \phm{}?
\item Apakah \asing{band} ketidakpastian mengurangi panggilan LLM dengan tetap menjaga mutu memori dan akurasi jawaban pada sesi lanjutan?
\end{enumerate}

\subbab{1.3 Tujuan dan Manfaat Penelitian}{sec:1.3}

\subsubbab{1.3.1 Tujuan Penelitian}{sec:1.3.1}

Penelitian ini bertujuan mengimplementasikan \asing{gate} klasifikasi terkalibrasi pada jalur tulis memori Pi dengan memisahkan penilaian kelayakan dan penulisan entri. Tujuan evaluasinya adalah sebagai berikut.
\begin{enumerate}[leftmargin=1.25cm]
\item Membandingkan akurasi keputusan \asing{gate} dengan mekanisme penyimpanan pada \pim{} dan \phm{} menggunakan \asing{dataset} berlabel dari sesi \asing{coding agent}.
\item Menilai kalibrasi probabilitas \asing{gate} dan kelayakannya sebagai dasar penentuan ambang keputusan.
\item Membandingkan penggunaan token, biaya, dan latensi per operasi tulis antara mekanisme yang diusulkan, \pim{}, dan \phm{}.
\item Menguji pengaruh \asing{band} ketidakpastian terhadap jumlah panggilan LLM, mutu memori, dan akurasi jawaban pada sesi lanjutan.
\end{enumerate}

\subsubbab{1.3.2 Manfaat Penelitian}{sec:1.3.2}

\textbf{Manfaat teoretis.} Penelitian ini menambah pemahaman tentang penerapan \asing{gate} klasifikasi terkalibrasi dalam menilai kelayakan penyimpanan memori sesi \asing{coding agent}.

\textbf{Manfaat praktis.} Hasil penelitian dapat membantu pengguna ekosistem Pi memilih mekanisme penyimpanan memori berdasarkan akurasi keputusan, penggunaan token, biaya, dan latensi per operasi tulis.

\textbf{Manfaat metodologis.} Penelitian ini menyediakan protokol evaluasi \asing{gate} memori yang mencakup data berlabel, pemeriksaan konsistensi pelabelan, kalibrasi, dan pengujian mutu memori pada sesi lanjutan.

% ------------------------------------------------------------------
% BAB II
% ------------------------------------------------------------------
\bab{II}{KAJIAN PUSTAKA DAN LANDASAN TEORI}{sec:bab2}

\subbab{2.1 Kajian Pustaka dan Landasan Teori}{sec:2.1}

\subsubbab{2.1.1 Memori Lintas Sesi pada \asing{Coding Agent}}{sec:2.1.1}

Memori lintas sesi menyimpan informasi yang dapat membantu agen pada pekerjaan berikutnya, seperti preferensi pengguna, konvensi repositori, keputusan arsitektur, dan pelajaran dari kegagalan \asing{tool}. Dokumentasi \pim{} menjelaskan penyimpanan melalui berkas \asing{markdown}, log harian, \asing{scratchpad}, serta \asing{tool} \texttt{memory\_write} (jayzeng, 2026). Sementara itu, \phm{} menggunakan memori berkategori dan \asing{review} percakapan untuk menghasilkan pembaruan memori (chandra447, 2026).

Penelitian ini berfokus pada memori jangka panjang. Informasi dinilai layak disimpan jika mempunyai dasar dalam percakapan atau hasil pekerjaan, sesuai dengan cakupan pengguna atau proyek, dan berguna untuk sesi berikutnya. Informasi sementara tetap dapat digunakan dalam konteks aktif atau log.

\subsubbab{2.1.2 Penilaian Kelayakan dan Kendali Jalur Tulis Memori}{sec:2.1.2}

Kendali jalur tulis memori atau \asing{write-side control} menentukan tindakan terhadap informasi sebelum disimpan. SAGE menilai apakah informasi sudah tercakup dalam memori melalui estimator kepadatan berbasis distribusi von Mises--Fisher pada \asing{embedding}. Informasi yang jelas belum tercakup ditambahkan, informasi berulang diabaikan, dan kasus yang belum jelas diperiksa oleh LLM (Wang et al., 2026).

Informasi yang belum tersimpan belum tentu berguna untuk jangka panjang. Koreksi terhadap keputusan lama juga dapat tetap penting meskipun teksnya mirip dengan memori sebelumnya. Karena itu, penilaian kelayakan mempertimbangkan kegunaan informasi dan konteksnya. Command Code Taste juga memisahkan penilaian preferensi dari penulisan hasil oleh LLM, tetapi penghematan pada Pi tetap perlu diuji (Command Code, 2026).

\subsubbab{2.1.3 Klasifikasi Kelayakan Penyimpanan}{sec:2.1.3}

Kelayakan penyimpanan dinilai melalui klasifikasi biner. Kelas positif menunjukkan informasi yang layak masuk ke memori jangka panjang, sedangkan kelas negatif menunjukkan informasi yang tidak perlu disimpan di sana. Label acuan disusun melalui pedoman anotasi, dan konsistensi pelabelan diperiksa sebelum label digunakan. Ketepatan \asing{gate} diukur menggunakan akurasi, \asing{precision}, \asing{recall}, dan F1 agar kesalahan menerima maupun menolak informasi dapat diketahui.

\subsubbab{2.1.4 Kalibrasi Probabilitas \asing{Gate}}{sec:2.1.4}

Kalibrasi menunjukkan kesesuaian probabilitas prediksi dengan hasil yang diamati. Model dengan akurasi tinggi belum tentu menghasilkan probabilitas yang terkalibrasi (Guo et al., 2017). Karena itu, ketepatan kelas dan mutu probabilitas dinilai secara terpisah.

Guo et al. (2017) menggunakan \asing{reliability diagram}, \asing{Expected Calibration Error} (ECE), dan \asing{negative log-likelihood} (NLL) untuk mengevaluasi kalibrasi. ECE mengukur ketidaksesuaian probabilitas dengan hasil pengamatan. Salah satu metode yang dibahas adalah \asing{temperature scaling}, dengan parameter yang disesuaikan menggunakan data validasi. Metode ini mengubah probabilitas tanpa mengubah kelas yang dipilih, sehingga perbaikan kalibrasi tidak otomatis meningkatkan akurasi klasifikasi.

\subsubbab{2.1.5 \asing{Band} Ketidakpastian dan Pemeriksaan LLM}{sec:2.1.5}

\asing{Band} ketidakpastian digunakan untuk meneruskan kandidat ke LLM ketika keluaran \asing{gate} belum cukup meyakinkan. SAGE mengalihkan kasus yang belum jelas berdasarkan penilaian apakah informasi sudah tercakup dalam memori (Wang et al., 2026). Dalam penelitian ini, pengalihan direncanakan berdasarkan probabilitas kelayakan penyimpanan. Penentuan ambang dan rentangnya dijelaskan pada bagian metode penelitian. Pengurangan panggilan LLM dinilai bersama kesalahan keputusan dan mutu memori.

\subsubbab{2.1.6 Evaluasi Efisiensi dan Mutu Memori}{sec:2.1.6}

Efisiensi dinilai melalui penggunaan token, biaya, dan latensi. Pengukuran mencakup ekstraksi kandidat, komputasi \asing{gate}, pemeriksaan LLM, serta penulisan entri. Biaya untuk kandidat yang akhirnya ditolak tetap dihitung. Perbandingan menggunakan definisi operasi tulis yang sama dan mencatat perbedaan konfigurasi kedua \asing{baseline}, termasuk \asing{review} berkala pada \phm{} (jayzeng, 2026; chandra447, 2026).

Maharana et al. (2024) mengevaluasi penggunaan memori percakapan melalui beberapa tugas, termasuk \asing{question answering}. Dalam penelitian ini, mutu memori diuji melalui ketepatan jawaban sesi lanjutan tentang keputusan teknis, konvensi repositori, dan koreksi pengguna yang tersedia dalam histori.

\subsubbab{2.1.7 Posisi Penelitian}{sec:2.1.7}

SAGE telah membahas pemisahan sebagian keputusan penyimpanan dari pemeriksaan LLM (Wang et al., 2026). Penelitian ini berfokus pada evaluasi komparatif \asing{gate} klasifikasi terkalibrasi pada \asing{coding agent} Pi, dengan menilai ketepatan keputusan, mutu probabilitas, efisiensi, dan penggunaan memori pada sesi lanjutan.

\subbab{2.2 Pertanyaan Penelitian dan Hipotesis}{sec:2.2}

Pertanyaan penelitian mengikuti rumusan masalah pada BAB I. Hipotesis berikut merupakan dugaan yang akan diuji menggunakan data yang tidak digunakan untuk pelatihan atau penyesuaian kalibrasi.

\textbf{Pertanyaan penelitian 1.} Bagaimana akurasi keputusan \asing{gate} klasifikasi terkalibrasi dibandingkan dengan mekanisme penyimpanan pada \pim{} dan \phm{}?

\textbf{Hipotesis 1.} Terdapat perbedaan ketepatan keputusan kelayakan penyimpanan antara \asing{gate} klasifikasi terkalibrasi dan kedua \asing{baseline}. Perbandingan menggunakan label acuan yang sama melalui akurasi, \asing{precision}, \asing{recall}, dan F1.

\textbf{Pertanyaan penelitian 2.} Apakah probabilitas yang dihasilkan \asing{gate} terkalibrasi dan dapat digunakan untuk menentukan keputusan?

\textbf{Hipotesis 2.} Kalibrasi \asing{post-hoc} membuat probabilitas \asing{gate} lebih sesuai dengan hasil pada data acuan, ditunjukkan oleh penurunan ECE dibandingkan sebelum kalibrasi. Kegunaan probabilitas tersebut untuk menentukan keputusan otomatis dinilai bersama tingkat kesalahan \asing{gate}.

\textbf{Pertanyaan penelitian 3.} Bagaimana perbandingan penggunaan token, biaya, dan latensi per operasi tulis antara mekanisme yang diusulkan, \pim{}, dan \phm{}?

\textbf{Hipotesis 3.} \asing{Gate} dengan pemeriksaan LLM selektif menurunkan penggunaan token LLM dan biaya operasional per operasi tulis dibandingkan \asing{baseline}, setelah memperhitungkan komputasi tambahan \asing{gate}. Latensi dinilai tersendiri karena pengurangan panggilan LLM belum tentu mempercepat seluruh operasi.

\textbf{Pertanyaan penelitian 4.} Apakah \asing{band} ketidakpastian mengurangi panggilan LLM dengan tetap menjaga mutu memori dan akurasi jawaban pada sesi lanjutan?

\textbf{Hipotesis 4.} \asing{Band} ketidakpastian mengurangi panggilan LLM untuk penilaian kelayakan, sementara mutu memori dan akurasi jawaban sesi lanjutan tidak menurun melebihi batas yang ditetapkan sebelum pengujian. Dengan demikian, penghematan dinilai bersama kemampuan agen menggunakan informasi yang diperlukan.

% ------------------------------------------------------------------
% BAB III
% ------------------------------------------------------------------
\bab{III}{METODE PENELITIAN}{sec:bab3}

\subbab{3.1 Metode Penelitian}{sec:3.1}

\subsubbab{3.1.1 Pendekatan dan Jenis Penelitian}{sec:3.1.1}

Penelitian ini akan menggunakan pendekatan kuantitatif dengan metode eksperimen komparatif. Pendekatan kuantitatif dipilih karena seluruh pertanyaan penelitian menanyakan besaran yang dapat diukur, yaitu ketepatan keputusan, kalibrasi, penggunaan token, biaya, latensi, dan akurasi jawaban pada sesi lanjutan. Metode eksperimen dipilih karena mekanisme penyimpanan dapat dikendalikan sementara kondisi lain dibuat sama, sehingga perbedaan hasil dapat dikaitkan dengan mekanisme yang digunakan. Survei dan studi kasus tidak dipilih karena penelitian ini mengukur perilaku sistem pada data yang sama.

\subsubbab{3.1.2 Objek, Populasi, dan Sampel Penelitian}{sec:3.1.2}

Objek penelitian adalah mekanisme penilaian kelayakan informasi sebelum disimpan ke memori jangka panjang \asing{coding agent} Pi. Unit analisisnya adalah satu kandidat memori beserta konteks percakapan dan isi memori yang tersedia saat kandidat tersebut muncul.

Populasi penelitian adalah seluruh kandidat memori dari sesi \asing{coding agent} Pi milik peneliti selama periode pengumpulan data. Sampel diambil dari sesi yang memuat pekerjaan pemrograman nyata, seperti perbaikan kode, perubahan konfigurasi, dan diskusi keputusan teknis. Sampel dibagi per sesi menjadi data latih, data validasi, dan data uji, sehingga kandidat dari sesi yang sama tidak muncul pada kelompok yang berbeda.

\subsubbab{3.1.3 Variabel Penelitian}{sec:3.1.3}

Variabel bebas penelitian ini adalah mekanisme penilaian kelayakan penyimpanan, yaitu mekanisme bawaan \pim{}, mekanisme bawaan \phm{}, dan \asing{gate} klasifikasi terkalibrasi. Untuk pertanyaan penelitian keempat, lebar \asing{band} ketidakpastian juga diperlakukan sebagai variabel bebas. Variabel terikat meliputi ketepatan keputusan (akurasi, \asing{precision}, \asing{recall}, dan F1), mutu probabilitas \asing{gate} (ECE dan NLL) (Guo et al., 2017), efisiensi (jumlah panggilan LLM, token, biaya, dan latensi per operasi tulis), serta mutu memori (akurasi jawaban pada sesi lanjutan). Variabel kontrol meliputi model LLM dan versinya, instruksi penulisan entri, kondisi awal memori, dan perangkat keras yang digunakan.

\subsubbab{3.1.4 Desain dan Prosedur Eksperimen}{sec:3.1.4}

Desain eksperimen yang digunakan adalah perbandingan dalam subjek (\asing{within-subject}), yaitu setiap kandidat memori pada data uji diproses oleh ketiga mekanisme. Penelitian ini tidak menguji arsitektur memori agen secara keseluruhan seperti pada MemGPT atau A-MEM (Packer et al., 2023; Xu et al., 2025). Pengujian dibatasi pada tahap keputusan sebelum informasi ditulis ke memori. Prosedur eksperimen direncanakan sebagai berikut.
\begin{enumerate}[leftmargin=1.25cm]
\item Mengumpulkan rekaman sesi Pi, menyamarkan data pribadi, lalu mengekstraksi kandidat memori dengan aturan yang sama untuk ketiga mekanisme.
\item Memberi label acuan berdasarkan pedoman anotasi, lalu membagi data per sesi menjadi data latih, validasi, dan uji.
\item Melatih \asing{gate} pada data latih dan menyesuaikan kalibrasi dengan \asing{temperature scaling} pada data validasi (Guo et al., 2017).
\item Menentukan ambang keputusan dan lebar \asing{band} ketidakpastian pada data validasi. Kandidat di dalam \asing{band} diteruskan ke LLM, mengikuti prinsip klasifikasi selektif yang menahan keputusan pada prediksi berkeyakinan rendah (Geifman \& El-Yaniv, 2017).
\item Memutar ulang sesi data uji pada ketiga mekanisme dengan konfigurasi yang sama, lalu mencatat keputusan, jumlah panggilan LLM, token, biaya, dan latensi setiap kandidat. Pada \pim{} dan \phm{}, kandidat dianggap diterima apabila informasinya tertulis ke memori.
\item Menjalankan sesi lanjutan berisi pertanyaan tentang keputusan teknis, konvensi repositori, dan koreksi pengguna, mengikuti kategori evaluasi memori jangka panjang pada Maharana et al. (2024) dan Wu et al. (2025), lalu mencatat ketepatan jawaban agen.
\end{enumerate}

Ambang, kalibrasi, dan lebar \asing{band} ditetapkan sebelum data uji digunakan.

\subbab{3.2 Teknik Pengumpulan Data dan Instrumen Penelitian}{sec:3.2}

\subsubbab{3.2.1 Teknik Pengumpulan Data}{sec:3.2.1}

Data akan dikumpulkan melalui dokumentasi rekaman sesi, anotasi kandidat memori oleh peneliti, dan pencatatan otomatis (\asing{logging}) selama eksperimen. Nama, alamat surel, kredensial, dan data sensitif lain dalam rekaman dihapus atau disamarkan sebelum digunakan. Kesesuaian teknik dengan setiap variabel ditunjukkan pada Tabel 3.1.

\begin{table}[h]
\centering
\captionsetup{labelformat=empty,font=small,skip=4pt}
\caption{\textbf{Tabel 3.1} Kesesuaian variabel, teknik pengumpulan data, dan instrumen}
\begin{singlespace}
\small
\begin{tabularx}{\linewidth}{>{\raggedright\arraybackslash}X >{\raggedright\arraybackslash}X >{\raggedright\arraybackslash}X}
\toprule
\textbf{Variabel} & \textbf{Teknik pengumpulan data} & \textbf{Instrumen} \\
\midrule
Mekanisme penyimpanan (variabel bebas) & Konfigurasi eksperimen & Berkas konfigurasi dan catatan versi kode \\
Ketepatan keputusan & Anotasi dan pencatatan otomatis & Pedoman anotasi, lembar pelabelan, log keputusan \\
Mutu probabilitas \asing{gate} & Pencatatan otomatis & Skrip pengujian dan log probabilitas \\
Efisiensi & Pencatatan otomatis & Log panggilan LLM, token, biaya, dan waktu proses \\
Mutu memori & Dokumentasi dan pencatatan otomatis & Daftar pertanyaan dan jawaban acuan sesi lanjutan \\
\bottomrule
\end{tabularx}
\end{singlespace}
\end{table}

\subsubbab{3.2.2 Instrumen Penelitian}{sec:3.2.2}

Instrumen penelitian terdiri atas empat bagian.
\begin{enumerate}[leftmargin=1.25cm]
\item \textbf{Pedoman anotasi}, berisi definisi kelas ``layak disimpan'' dan ``tidak layak disimpan'' beserta contohnya. Preferensi pengguna yang relatif menetap, konvensi repositori, dan keputusan teknis yang masih berlaku termasuk layak disimpan. Informasi sementara, duplikasi, dan klaim tanpa dasar termasuk tidak layak disimpan. Koreksi terhadap memori lama diperlakukan sebagai informasi baru meskipun teksnya mirip.
\item \textbf{Lembar pelabelan}, mencatat identitas kandidat, konteks singkat, label, dan alasan pemberian label.
\item \textbf{Skrip pengujian dan pencatat log}, menjalankan ketiga mekanisme dengan konfigurasi yang sama dan mencatat keputusan, probabilitas \asing{gate}, panggilan LLM, token, biaya, serta waktu proses.
\item \textbf{Daftar pertanyaan sesi lanjutan}, setiap pertanyaan dikaitkan dengan kandidat berlabel ``layak disimpan'' dan dilengkapi jawaban acuan.
\end{enumerate}

\subsubbab{3.2.3 Validitas dan Reliabilitas Instrumen}{sec:3.2.3}

Validitas pedoman anotasi dijaga dengan menurunkan definisi kelas dari kriteria kelayakan pada subbab 2.1.1, yaitu informasi memiliki dasar dalam percakapan atau hasil pekerjaan, sesuai dengan cakupan pengguna atau proyek, dan berguna untuk sesi berikutnya. Pedoman diujicobakan pada sebagian kecil kandidat dan diperbaiki apabila terdapat definisi yang menimbulkan keraguan. Setiap pertanyaan sesi lanjutan disusun agar hanya dapat dijawab dengan benar apabila informasi terkait tersimpan di memori.

Pelabelan dilakukan oleh satu penilai, yaitu peneliti. Karena itu, reliabilitas diperiksa melalui pelabelan ulang. Peneliti melabeli kembali sebagian kandidat setelah jeda waktu tanpa melihat label pertama, lalu kesepakatan kedua pelabelan dihitung dengan koefisien kappa Cohen (Cohen, 1960). Penggunaan satu penilai dicatat sebagai keterbatasan penelitian. Reliabilitas pengukuran efisiensi dijaga dengan mengulang pengukuran latensi pada konfigurasi yang sama serta mencatat versi model, instruksi, dan kode yang digunakan.

% ------------------------------------------------------------------
% BAB IV
% ------------------------------------------------------------------
\bab{IV}{ANALISIS DATA DAN KESIMPULAN SEMENTARA}{sec:bab4}

\subbab{4.1 Teknik Analisis Data Penelitian}{sec:4.1}

Data penelitian ini berupa angka, yaitu hasil keputusan setiap mekanisme, probabilitas \asing{gate}, penggunaan token, biaya, latensi, dan ketepatan jawaban pada sesi lanjutan. Karena itu, data akan dianalisis dengan statistika deskriptif dan statistika inferensial (Santosa, 2021). Pengolahan data akan dilakukan menggunakan bahasa pemrograman Python.

\subsubbab{4.1.1 Statistika Deskriptif}{sec:4.1.1}

Statistika deskriptif digunakan untuk meringkas data agar gambaran awal hasil setiap mekanisme dapat diketahui (Santosa, 2021). Sebelum perhitungan dilakukan, kesepakatan antara pelabelan pertama dan pelabelan ulang diperiksa dengan koefisien kappa Cohen seperti dijelaskan pada subbab 3.2.3. Setelah itu, untuk setiap mekanisme akan dihitung hal-hal berikut.
\begin{enumerate}[leftmargin=1.25cm]
\item Ketepatan keputusan, yaitu akurasi, \asing{precision}, \asing{recall}, dan F1. \asing{Precision} menunjukkan seberapa sedikit informasi yang salah disimpan, \asing{recall} menunjukkan seberapa sedikit informasi penting yang terlewat, dan F1 merupakan gabungan keduanya.
\item Nilai ECE \asing{gate} sebelum dan sesudah kalibrasi. Nilai ECE yang lebih kecil menunjukkan tingkat keyakinan \asing{gate} lebih sesuai dengan hasil sebenarnya.
\item Rerata, nilai tengah, dan simpang baku penggunaan token, biaya, dan latensi per operasi tulis. Biaya dihitung dari jumlah token dikali tarif normal tanpa potongan \asing{prompt caching}, sehingga ketiga mekanisme dinilai dengan aturan biaya yang sama.
\item Jumlah panggilan LLM dan akurasi jawaban sesi lanjutan pada kondisi dengan dan tanpa \asing{band} ketidakpastian.
\end{enumerate}
Hasil perhitungan akan disajikan dalam tabel per mekanisme agar perbandingan antarmekanisme mudah dibaca.

\subsubbab{4.1.2 Statistika Inferensial dan Uji Hipotesis}{sec:4.1.2}

Statistika inferensial digunakan untuk menguji hipotesis (Santosa, 2021). Hipotesis penelitian pada subbab 2.2 akan diubah menjadi hipotesis statistik, yaitu hipotesis nol (H\textsubscript{0}) yang menyatakan tidak ada perbedaan dan hipotesis alternatif (H\textsubscript{a}) yang menyatakan ada perbedaan. Pengujian akan menggunakan uji t dengan taraf signifikansi 5\%. Artinya, perbedaan hasil dianggap nyata apabila peluang perbedaan tersebut muncul karena kebetulan kurang dari 5\% (nilai p \textless{} 0,05). Apabila data tidak berdistribusi normal, seperti data latensi yang dapat memiliki beberapa nilai sangat besar, akan digunakan prosedur statistika nonparametrik (Santosa, 2021). Pengujian setiap hipotesis direncanakan sebagai berikut.
\begin{enumerate}[leftmargin=1.25cm]
\item \textbf{Hipotesis 1.} Nilai F1 setiap sesi pada data uji dibandingkan antara \asing{gate} terkalibrasi dan setiap \asing{baseline} menggunakan uji t.
\item \textbf{Hipotesis 2.} Nilai ECE sebelum dan sesudah kalibrasi dibandingkan. Hipotesis didukung apabila nilai ECE sesudah kalibrasi lebih kecil.
\item \textbf{Hipotesis 3.} Penggunaan token dan biaya per operasi tulis dibandingkan antara \asing{gate} terkalibrasi dan setiap \asing{baseline} menggunakan uji t. Latensi diuji dengan cara yang sama dan dilaporkan secara terpisah.
\item \textbf{Hipotesis 4.} Pengurangan panggilan LLM dihitung dalam bentuk persentase. Akurasi jawaban sesi lanjutan dianggap terjaga apabila penurunannya tidak lebih dari 5 poin persentase dibandingkan kondisi tanpa \asing{band} ketidakpastian.
\end{enumerate}

\subbab{4.2 Kesimpulan Sementara dan Batasan Penelitian}{sec:4.2}

\subsubbab{4.2.1 Kesimpulan Sementara}{sec:4.2.1}

Berdasarkan rumusan masalah dan hipotesis penelitian, penelitian ini diharapkan menunjukkan hal-hal berikut.
\begin{enumerate}[leftmargin=1.25cm]
\item Ketepatan keputusan \asing{gate} klasifikasi terkalibrasi diperkirakan berbeda dengan \pim{} dan \phm{}.
\item Kalibrasi diperkirakan menurunkan nilai ECE, sehingga probabilitas \asing{gate} dapat digunakan sebagai dasar penentuan ambang keputusan.
\item \asing{Gate} dengan pemeriksaan LLM selektif diperkirakan menggunakan token dan biaya yang lebih rendah per operasi tulis, sedangkan latensinya belum tentu lebih rendah.
\item \asing{Band} ketidakpastian diperkirakan mengurangi panggilan LLM tanpa menurunkan akurasi jawaban sesi lanjutan melebihi batas yang ditetapkan.
\end{enumerate}
Apabila dugaan tersebut terbukti, penelitian ini diharapkan dapat menjadi dasar bagi pengguna Pi dalam memilih mekanisme penyimpanan memori yang lebih hemat dengan mutu memori yang tetap terjaga. Dugaan ini akan diterima atau ditolak berdasarkan hasil pengujian pada data uji.

\subsubbab{4.2.2 Batasan Penelitian}{sec:4.2.2}

Keterbatasan penelitian adalah karakteristik rancangan atau metodologi penelitian yang dapat memengaruhi penafsiran temuan penelitian (Santosa, 2021). Keterbatasan penelitian ini beserta saran untuk penelitian berikutnya adalah sebagai berikut.
\begin{enumerate}[leftmargin=1.25cm]
\item \textbf{Proses pelabelan.} Pelabelan dilakukan oleh satu penilai, yaitu peneliti, sehingga label dapat dipengaruhi penilaian pribadi. Penelitian berikutnya dapat melibatkan penilai kedua.
\item \textbf{Profil sampel.} Data hanya berasal dari sesi \asing{coding agent} milik peneliti, sehingga hasilnya belum dapat diberlakukan untuk pengguna dan proyek lain. Penelitian berikutnya dapat menggunakan sesi dari beberapa pengguna.
\item \textbf{Ukuran sampel.} Jumlah sesi yang dapat dikumpulkan terbatas, sehingga hasil uji hipotesis perlu ditafsirkan dengan hati-hati. Penelitian berikutnya dapat menggunakan periode pengumpulan data yang lebih panjang.
\item \textbf{Ketergantungan pada model.} Hasil pengukuran bergantung pada versi model LLM dan harga API pada saat pengujian. Penelitian berikutnya dapat mengulang pengujian pada model lain.
\item \textbf{Cakupan rancangan.} Penelitian hanya menguji tahap keputusan sebelum informasi disimpan dan tidak menguji tahap pengambilan memori. Penelitian berikutnya dapat mengevaluasi tahap pengambilan memori.
\end{enumerate}

% ------------------------------------------------------------------
% DAFTAR PUSTAKA (APA Style edisi 7)
% ------------------------------------------------------------------
\clearpage
\phantomsection\label{sec:daftar-pustaka}
\begin{center}\large\bfseries DAFTAR PUSTAKA\end{center}
\vspace{0.8em}

\begin{singlespace}
\raggedright
\begin{list}{}{%
  \setlength{\leftmargin}{1.25cm}%
  \setlength{\itemindent}{-1.25cm}%
  \setlength{\itemsep}{0.8\baselineskip}%
  \setlength{\parsep}{0pt}%
}
\item chandra447. (2026). \textit{Pi Hermes Memory} [Dokumentasi perangkat lunak]. GitHub. Diakses pada 30 September 2026, dari \url{https://github.com/chandra447/pi-hermes-memory}

\item Cohen, J. (1960). A coefficient of agreement for nominal scales. \textit{Educational and Psychological Measurement, 20}(1), 37--46. \url{https://doi.org/10.1177/001316446002000104}

\item Command Code. (2026). \textit{Taste}. Command Code Docs. Diakses pada 30 September 2026, dari \url{https://commandcode.ai/docs/taste}

\item Geifman, Y., \& El-Yaniv, R. (2017). Selective classification for deep neural networks. \textit{Advances in Neural Information Processing Systems, 30}. \url{https://arxiv.org/abs/1705.08500}

\item Guo, C., Pleiss, G., Sun, Y., \& Weinberger, K. Q. (2017). On calibration of modern neural networks. Dalam D. Precup \& Y. W. Teh (Ed.), \textit{Proceedings of the 34th International Conference on Machine Learning} (Vol. 70, hlm. 1321--1330). PMLR. \url{https://proceedings.mlr.press/v70/guo17a.html}

\item jayzeng. (2026). \textit{pi-memory} [Dokumentasi perangkat lunak]. GitHub. Diakses pada 30 September 2026, dari \url{https://github.com/jayzeng/pi-memory}

\item Maharana, A., Lee, D.-H., Tulyakov, S., Bansal, M., Barbieri, F., \& Fang, Y. (2024). \textit{Evaluating very long-term conversational memory of LLM agents} (arXiv:2402.17753) [Preprint]. arXiv. \url{https://arxiv.org/abs/2402.17753}

\item Packer, C., Wooders, S., Lin, K., Fang, V., Patil, S. G., Stoica, I., \& Gonzalez, J. E. (2023). \textit{MemGPT: Towards LLMs as operating systems} (arXiv:2310.08560) [Preprint]. arXiv. \url{https://arxiv.org/abs/2310.08560}

\item Santosa, P. I. (2021). \textit{Metodologi penelitian} (Edisi 1). Universitas Terbuka.
\item Wang, S., Brahma, D., \& Henao, R. (2026). \textit{SAGE: A novelty gate for efficient memory evolution in agentic LLMs} (arXiv:2605.30711) [Preprint]. arXiv. \url{https://arxiv.org/abs/2605.30711v2}

\item Wu, D., Wang, H., Yu, W., Zhang, Y., Chang, K.-W., \& Yu, D. (2025). LongMemEval: Benchmarking chat assistants on long-term interactive memory. \textit{International Conference on Learning Representations (ICLR 2025)}. \url{https://arxiv.org/abs/2410.10813}

\item Xu, W., Liang, Z., Mei, K., Gao, H., Tan, J., \& Zhang, Y. (2025). \textit{A-MEM: Agentic memory for LLM agents} (arXiv:2502.12110) [Preprint]. arXiv. \url{https://arxiv.org/abs/2502.12110}
\end{list}
\end{singlespace}

\end{document}
